\begin{figure}[H]
\centering
\resizebox{0.70\textwidth}{!}{%
\begin{tikzpicture}[
  font=\sffamily,
  n/.style={circle, draw=gbnavy, fill=white, line width=.8pt, minimum size=5.5mm, inner sep=0pt},
  e/.style={-{Latex[length=1.8mm]}, draw=gbgray, line width=.7pt},
  panelcaption/.style={font=\sffamily\footnotesize, align=center}
]
% N1
\node[n] (a1) at (0,0) {1}; \node[n] (a2) at (1.3,0) {2}; \node[n] (a3) at (2.6,0) {3};
\draw[e](a1)--(a2); \draw[e](a2)--(a3);
\node[panelcaption] (capa) at (1.3,-.75) {N1: linear\\depth/attenuation};
% N2
\node[n] (b1) at (4.6,0) {1}; \node[n] (b2) at (6.0,.45) {2}; \node[n] (b3) at (6.0,-.45) {3};
\draw[e](b1)--(b2); \draw[e](b1)--(b3);
\node[panelcaption] (capb) at (5.3,-1.2) {N2: branching\\divergence};
% N3
\node[n] (c1) at (8.0,.45) {1}; \node[n] (c2) at (8.0,-.45) {2}; \node[n] (c3) at (9.5,0) {3};
\draw[e](c1)--(c3); \draw[e](c2)--(c3);
\node[panelcaption] (capc) at (8.75,-1.2) {N3: converging\\multi-source aggregation};
% N4
\node[n] (d1) at (1.0,-3.2) {1}; \node[n] (d2) at (2.3,-2.7) {2}; \node[n] (d3) at (2.3,-3.7) {3}; \node[n] (d4) at (3.7,-3.2) {4};
\draw[e](d1)--(d2); \draw[e](d1)--(d3); \draw[e](d2)--(d4); \draw[e](d3)--(d4);
\node[panelcaption] (capd) at (2.35,-4.45) {N4: reconverging\\shared-origin dependence};
% N5
\node[n] (m1) at (6.2,-2.7) {1}; \node[n] (m2) at (6.2,-3.7) {2}; \node[n] (m3) at (7.6,-2.5) {3}; \node[n] (m4) at (7.6,-3.5) {4}; \node[n] (m5) at (9.1,-3.1) {5}; \node[n] (m6) at (10.3,-3.1) {6};
\draw[e](m1)--(m3); \draw[e](m1)--(m4); \draw[e](m2)--(m4); \draw[e](m3)--(m5); \draw[e](m4)--(m5); \draw[e](m5)--(m6);
\node[panelcaption] (cape) at (8.25,-4.45) {N5: mixed\\combined mechanisms};
\end{tikzpicture}}
\caption{Controlled synthetic DAG set. Node labels identify evaluation order; arrows indicate permitted propagation direction.}
\label{fig:network-set}
\end{figure}
